\uuid{frG6}
\titre{Loi sur trois valeurs à paramètre}
\auteur{worms}
\organisation{exo7}
\licence{CC BY-NC-SA 4.0}
\datecreate{2017-08-28}
\niveau{L2}
\module{Probabilité et statistique}
\chapitre{Probabilité discrète}
\sousChapitre{Variable aléatoire discrète}
\difficulte{}
\theme{Variables discrètes}
\exo7qcmid{worms.worms.034}
\serie{ordonnee}

\contenu{
\texte{
Soit $X$ une variable aléatoire à valeurs dans $\{0,1,2\}$ et de loi donnée  par
\[
 \mathbb{P}(X=0)=\mathbb{P}(X=2)=a \makebox[1.cm][c]{et} \mathbb{P}(X=1)=1-2a
\]
où $a$ est une constante réelle.
}

\begin{question}
\qclassification{260.03}
\qkeeporder
Quelles valeurs la constante $a$ a-t-elle le droit de prendre ?
\begin{answers}
\bad{Toutes les valeurs de $]0,1[$ car $\mathbb{P}(X=0)+\mathbb{P}(X=1)+\mathbb{P}(X=2)=1$.}
\bad{Seulement la valeur $a=1/4$.}
\good{Toutes les valeurs de $]0,1/2[$.}
\bad{Une autre réponse que les précédentes.}
\end{answers}
\begin{explanations} Les probabilités $\mathbb{P}(X=k)$ doivent appartenir à $]0,1[$ (ni $0$ ni $1$, sinon une ou plusieurs modalités ne pourraient pas être déclarées dans l'espace d'état de $X$), d'où la réponse.
\end{explanations}
\end{question}

\begin{question}
\qclassification{260.03}
Quel est le graphe de la fonction de répartition de $X$ parmi les graphes suivants ?

%[[image]]
\qimage{images/png/frG6-1.png}
%
%\begin{center}
%\includegraphics[width=14.cm,height=5.cm]{ImagePourQCM2}
%\end{center}

\begin{answers}
\bad{Le premier.}
\good{Le second.}
\bad{Le troisième.}
\end{answers}
\end{question}

\begin{question}
\qclassification{260.03}
Que valent l'espérance et la variance de $X$ ?
\begin{answers}
\bad{$\mathbb{E}(X)=1$ et $\operatorname{Var}(X)=1+2a$.}
\bad{$\mathbb{E}(X)=2a$ et $\operatorname{Var}(X)=4a^2$.}
\good{$\mathbb{E}(X)=1$ et $\operatorname{Var}(X)=2a$.}
\end{answers}
\end{question}

\begin{question}
\qclassification{260.03}
\qkeeporder
On pose $Y=4-2X$. \emph{Sans déterminer la loi de $Y$}, peut-on calculer l'espérance et l'écart-type de $Y$ ?
\begin{answers}
\good{Oui, ils valent respectivement $2$ et $\sqrt{8a}$.}
\bad{Oui, ils valent respectivement $2$ et $\sqrt{4(1-a)}$.}
\bad{Oui, ils valent respectivement $4(1-a)$ et $4a$.}
\bad{Oui, mais aucune des propositions précédentes n'est correcte.}
\bad{Non, il nous faut nécessairement la loi pour calculer ces caractéristiques de $Y$.}
\end{answers}
\begin{explanations} $\mathbb{E}(Y)=4-2\mathbb{E}(X)$, et $\mathbb{V}ar(Y)=(-2)^2\mathbb{V}ar(X)$, donc l'écart-type de $Y$ vaut 2 fois celui de $X$.
\end{explanations}
\end{question}
}
